\section{Calibration-marginal adaptive testing}
\label{sec:framework}

\subsection{Items, oracles and recorded verdicts}
A \emph{test item} $i$ is a task together with an \emph{oracle} $O_i$ that maps a system
output to a verdict $y_i\in\{0,1\}$: exact numeric match for math word problems,
execution of hidden unit tests for code and repository-level software engineering, or a
log-likelihood comparison of fixed answer options for multiple-choice tasks. The framework
(called EOST in our project, for entropy-oracled information-optimal self-testing, where
optimality is meant in the one-step sense of Section~\ref{sec:theory}) imposes
one structural rule: \emph{the system under evaluation may propose or select items, but it
is never the oracle of its own verdicts}. In this paper all items come from fixed banks,
which isolates the measurement question; self-generated items were explored only in a
historical pilot (Section~\ref{sec:pilot}).

Two kinds of verdicts occur in our data and are kept distinct. A \emph{deterministic
verdict} is the output of the oracle on one recorded system output (GSM8K, ARC-C,
HellaSwag, WinoGrande, MMLU: a single greedy or log-likelihood evaluation; SWE-bench
Verified: the official harness on the submitted patch). A \emph{recorded aggregate
verdict} summarizes several sampled outputs: for LiveCodeBench and CRUXEval-O the public
data give a pass rate over samples, which we binarize at one half. Replaying either kind
is exact for the recorded data, but only the first corresponds to one administration of
the item to one deterministic system; Section~\ref{sec:study2} reports the main contrasts
separately for the two groups.

\subsection{Measurement model and calibration panel}
Let $s$ be the target system with latent competence $\theta_s$. Responses follow
the two-parameter logistic (2PL) model
\begin{equation}
\Pr(y_{si}=1\mid\theta_s,a_i,b_i)=\sigma\big(a_i(\theta_s-b_i)\big),\qquad a_i>0 .
\end{equation}
Item parameters $\psi_i=(\log a_i,b_i)$ are unknown and are calibrated on a
\emph{panel} $\mathcal C$ of $N$ other systems whose verdicts on the bank are
available. The panel excludes the target and every related system: in study~1 the
target's heuristic \emph{family} (same model line, reasoning-effort or version variant,
or, for agentic systems, any scaffold sharing an underlying LLM); in study~2 its
\emph{strict cluster}, which also joins near-duplicates detected from the verdicts
themselves (Section~\ref{sec:design}).

Calibration proceeds in four steps, shared by every policy:
(i) joint maximum a posteriori (MAP) estimation of the calibrators' abilities
$\hat\theta_{\mathcal C}$ and item parameters with $\theta_j\sim\mathcal N(0,1)$,
$\log a_i\sim\mathcal N(0,1)$, $b_i\sim\mathcal N(0,2^2)$;
(ii) an empirical-Bayes item hyperprior
$\log a_i\sim\mathcal N(\mu_a,\sigma_a^2)$, $b_i\sim\mathcal N(\mu_b,\sigma_b^2)$ obtained by
maximizing the marginal likelihood of \emph{all} items given
$\hat\theta_{\mathcal C}$ (EM on a parameter grid); the hyperparameters pool
hundreds of items and are therefore well determined even when $N=4$;
(iii) plug-in item parameters $(\hat a_i,\hat b_i)$, the conditional MAP given
$\hat\theta_{\mathcal C}$ under that hyperprior;
(iv) the exact grid posterior $\pi_i(\psi)\propto p(\psi)\prod_{j\in\mathcal C}p(y_{ji}\mid\hat\theta_j,\psi)$
on a $26\times129$ grid over $(\log a,b)\in[-2.5,2.5]\times[-8,8]$.

Two measurement models follow. The \textbf{plug-in} model uses
$P_i(\theta)=\sigma(\hat a_i(\theta-\hat b_i))$, as in classical computerized
adaptive testing and in recent LLM work \citep{hofmann2025fluid,li2026atlas}.
The \textbf{marginal} model uses the expected response function
\begin{equation}\label{eq:erf}
\bar P_i(\theta)=\int\sigma\big(a(\theta-b)\big)\,\pi_i(d\psi)
\end{equation}
\citep{mislevy1994erf}. By Lemma~\ref{lem:exact}, $\bar P_i$ is the exact response
probability of the target given $\hat\theta_{\mathcal C}$ and the hyperprior, so no
further approximation is involved in the ability posterior, the predictive distribution
or the information criteria below. Both models share $\hat\theta_{\mathcal C}$ and the
hyperprior: they differ only in plugging in a mode versus integrating the item
posterior.

Numerically, the joint MAP estimate of step (i) uses L-BFGS-B with analytic gradients. The
hyperprior of step (ii) is obtained by EM on a $14\times65$ grid (at most 60 iterations,
relative tolerance $10^{-7}$), with $\sigma_b\in[0.3,4]$ and $\sigma_a\in[0.1,1.5]$. The
plug-in parameters of step (iii) are computed by vectorized Fisher scoring with step halving.
For the grid posteriors of step (iv), with $\Lambda_1=\log\sigma(z)$ and
$\Lambda_0=\log\sigma(-z)$ evaluated on the $N\times G$ matrix $z_{jg}=a_g(\hat\theta_j-b_g)$,
the item log-likelihoods are $Y^\top\Lambda_1+(1-Y)^\top\Lambda_0$ (two GEMMs), and the
expected response functions on the ability grid are one further GEMM.

\subsection{Adaptive testing and selection criteria}
The target's ability posterior $\pi_t$ is kept on a 201-point grid over
$[-5,5]$ with prior $\mathcal N(0,1)$ and is updated after each verdict. At
step $t$ an item is chosen from the remaining candidate pool by one of:
\begin{itemize}\itemsep0pt
\item \textsc{Random}: uniform choice (the model is used only for scoring);
\item \textsc{Entropy}: maximize $U_{\mathrm H}(i)=h(\bar p_i)$, the predictive entropy
(uncertainty sampling);
\item \textsc{MFI}: maximize Fisher information at the posterior mean (classical CAT);
\item \textsc{PWFI}: maximize posterior-weighted Fisher information \citep{vanderlinden1998bayesian};
\item \textsc{EIG}: maximize $U_{\mathrm{EIG}}(i)=I(\theta;y_i\mid\mathcal D_t)$, the
expected information gain about competence (the EOST criterion);
\item \textsc{EVR}: maximize the expected reduction of the posterior variance of
the unobserved bank total (score-targeted EOST, Proposition~\ref{prop:evr}).
\end{itemize}
Each criterion is run under the plug-in (\texttt{:pi}) and the marginal
(\texttt{:mg}) model, except PWFI, which is run under the marginal model only. The proposed method, \textbf{EOST-M}, is \textsc{EIG:mg} (Algorithm~\ref{alg:eost}, Figure~\ref{fig:framework}).
The classical state of practice is \textsc{MFI:pi}. A model-free baseline,
\textsc{Naive}, samples items uniformly and reports the running sample mean
with a Wilson interval and finite-population correction.

The ability posterior is kept in log space. Entropy, EIG and PWFI are matrix--vector
products with precomputed item$\times$grid matrices; MFI interpolates Fisher information at
the posterior mean; EVR uses the closed form of Proposition~\ref{prop:evr} evaluated exactly
on the grid (unit-tested against a brute-force enumeration). Ties are broken uniformly at
random with a dedicated random stream.

\subsection{Recalibrating the score interval}
The model-based interval ignores what the model gets wrong (misspecification,
dependence between systems) and, as we will see, is too narrow on real data. A simple
correction uses only the calibration panel. For every calibrator $c$, the item posteriors
are recomputed without $c$'s cluster (abilities and hyperprior held at the panel fit), the
same adaptive test is replayed on $c$, whose full-bank score is known, and its absolute
score errors $|\hat A_{c,t}-A_c|$ are recorded. The recalibrated half-width at step $t$ is
the 90\% quantile of these errors pooled over calibrators and over the steps within a
factor 1.5 of $t$. This leave-one-calibrator-out rule (``+R'') is a cross-validated
correction, not a conformal guarantee: with $N\le8$ pseudo-targets no finite-sample
coverage guarantee at 90\% is available \citep{vovk2005algorithmic}. The rule was chosen
among three variants on development banks before study~2 (Section~\ref{sec:devhist}).

\subsection{Estimands and endpoints}
Latent competence is a model construct whose scale is fixed only by the prior,
so endpoints are defined on observable quantities.
\begin{enumerate}\itemsep0pt
\item \textbf{Held-out predictive log loss.} A fixed holdout $H$ (20\% of the
bank, drawn once per target and replicate and shared by \emph{all} policies and
panel sizes, a design adopted after an audit of the pilot program, Section~\ref{sec:pilot}) is never
administered. After each step, every policy predicts the target's verdicts on
$H$ with $\mathbb E_{\pi_t}P_i(\theta)$; the integrated log loss averages the
holdout log loss over steps $t=1,\dots,T$. A strictly proper score, it rewards
both accuracy and calibration of what the test has learned about the target.
\item \textbf{Score error.} The full-bank accuracy $A=\frac1I\sum_iy_{si}$ is
estimated by $\hat A_t=\frac1I\big(\sum_{\text{asked}}y_{si}+\sum_{\text{not asked}}\mathbb E_{\pi_t}P_i(\theta)\big)$;
the integrated absolute error (IAE) averages $|\hat A_t-A|$ over $t\le T$.
\item \textbf{Interval coverage and interval score.} The proportion of steps at which
the 90\% predictive interval for $A$ contains $A$, and the integrated interval score
$(u-l)+\frac2\alpha(l-A)_++\frac2\alpha(A-u)_+$ with $\alpha=0.1$, a proper score that
penalizes both width and misses \citep{gneiting2007strictly}. The predictive distribution of
the unobserved total is a mixture over the ability grid of sums of independent Bernoulli
variables; the interval uses its exact mean and variance and a normal approximation.
\item Secondary: holdout Brier score, final-step errors, sustained items-to-tolerance
(the first $t$ after which $|\hat A_s-A|\le\varepsilon$ for all $s\ge t$), and, for
synthetic banks, error and coverage of the true $\theta$.
\end{enumerate}
